\documentclass[11pt]{article}
\usepackage[margin=1in]{geometry}
\usepackage{amsmath,amssymb,amsthm}
\usepackage{algorithm}
\usepackage{algpseudocode}
\usepackage{booktabs}
\usepackage{hyperref}

\newcommand{\R}{\mathbb{R}}
\newcommand{\Ber}{\mathrm{Bernoulli}}
\newcommand{\logodds}{\ell}
\newcommand{\sigmoid}{\sigma}
\DeclareMathOperator{\softplus}{softplus}
\DeclareMathOperator{\logsumexp}{logsumexp}
\newtheorem{proposition}{Proposition}
\newtheorem{lemma}{Lemma}
\newtheorem{remark}{Remark}

\title{Incremental Path Verification:\\Approximate and Exact Bayesian Update Methods}
\author{}
\date{}

\begin{document}
\maketitle

% ====================================================================
\section{Problem Setup}
% ====================================================================

Consider a directed graph $G = (V, E)$ with $n = |V|$ vertices and
$m = |E|$ edges.  Each edge $e \in E$ has a hidden binary hazard state
$Z_e \in \{0, 1\}$.  The full hazard configuration is the random vector
$\mathbf{Z} = (Z_1, \dots, Z_m) \in \{0,1\}^m$.

\paragraph{Prior.}
Edges are a priori independent:
\[
  P(\mathbf{Z} = \mathbf{z})
  = \prod_{e=1}^{m} p_e^{z_e}\,(1-p_e)^{1-z_e},
  \qquad p_e = P(Z_e = 1) \in [0,1].
\]

\paragraph{Observation model.}
At each time step $t = 0, 1, 2, \dots$ we select a \emph{query mask}
$\mathbf{q}^{(t)} \in \{0,1\}^m$ indicating which edges are traversed,
and receive a single binary observation:
\[
  r^{(t)}
  = \bigvee_{e:\, q_e^{(t)} = 1} Z_e
  = \mathbf{1}\!\Bigl[\,\exists\, e : q_e^{(t)} = 1 \;\wedge\; Z_e = 1\Bigr]
  \;\in\;\{0,1\}.
\]
That is, $r = 1$ (\emph{collision}) if \emph{any} queried edge is
hazardous, and $r = 0$ (\emph{safe}) if \emph{all} queried edges are
hazard-free.  We denote the history of observations up to time $T$ by
\[
  \mathcal{D}_T = \bigl\{(\mathbf{q}^{(0)}, r^{(0)}),\;
                          (\mathbf{q}^{(1)}, r^{(1)}),\;
                          \dots,\;
                          (\mathbf{q}^{(T-1)}, r^{(T-1)})\bigr\}.
\]

\paragraph{Goal.}
Maintain a posterior belief $P(\mathbf{Z} \mid \mathcal{D}_t)$ and
quantify the \emph{realized information gain} at each step.

% ====================================================================
\section{Shared Definitions and Identities}
% ====================================================================

\subsection{Core Functions}

For a probability $p \in (0,1)$:
\begin{align}
  \logodds(p)     &= \log\frac{p}{1-p}       & &\text{(log-odds)} \label{eq:logodds} \\[4pt]
  \sigmoid(\ell)  &= \frac{1}{1+e^{-\ell}}   & &\text{(sigmoid, inverse of log-odds)} \label{eq:sigmoid} \\[4pt]
  \softplus(\ell) &= \log(1 + e^{\ell})       & &\text{(softplus)} \label{eq:softplus} \\[4pt]
  h(p)            &= -p\log_2 p - (1{-}p)\log_2(1{-}p) & &\text{(binary entropy, bits)} \label{eq:binent}
\end{align}
with boundary conventions $h(0) = h(1) = 0$, $\logodds(0) = -\infty$,
$\logodds(1) = +\infty$.

\subsection{Log-Probability Identities}

Given log-odds $\ell = \logodds(p)$:
\begin{align}
  \log p     &= \ell - \softplus(\ell)   = -\softplus(-\ell), \label{eq:logp} \\
  \log(1-p)  &= -\softplus(\ell). \label{eq:log1mp}
\end{align}
These follow from $p = e^\ell / (1 + e^\ell)$ and $1 - p = 1/(1 + e^\ell)$.

\subsection{Binary Entropy from Log-Odds}

The binary entropy $h(\sigmoid(\ell))$ can be computed directly from
$\ell$ without converting to probability.  Using
$\sigmoid(\ell) = e^\ell/(1+e^\ell)$ and $\sigmoid(-\ell) = 1 - \sigmoid(\ell)$:
\begin{align}
  h\bigl(\sigmoid(\ell)\bigr)
  &= -\sigmoid(\ell)\log_2\sigmoid(\ell) - \sigmoid(-\ell)\log_2\sigmoid(-\ell) \notag \\
  &= \frac{1}{\ln 2}\Bigl[\sigmoid(\ell)\,\softplus(-\ell) + \sigmoid(-\ell)\,\softplus(\ell)\Bigr].
  \label{eq:hlogodds}
\end{align}

\begin{proof}
Since $\log \sigmoid(\ell) = \ell - \log(1+e^\ell) = -\softplus(-\ell)$
and $\log \sigmoid(-\ell) = -\softplus(\ell)$, we have
$-\sigmoid(\ell)\log\sigmoid(\ell) = \sigmoid(\ell)\,\softplus(-\ell)$
and similarly for the second term.  Dividing by $\ln 2$ converts from
nats to bits.
\end{proof}

\subsection{Numerically Stable Helpers}

The implementation uses two additional stable primitives:

\paragraph{Log-sum-exp.}
For a vector $(x_1, \dots, x_n)$:
\[
  \logsumexp(x_1, \dots, x_n) = \log\!\sum_{i=1}^n e^{x_i}
  = x_{\max} + \log\!\sum_{i=1}^n e^{x_i - x_{\max}},
  \qquad x_{\max} = \max_i x_i.
\]

\paragraph{Log-one-minus-exp.}
For $a \leq 0$:
\[
  \texttt{log1mexp}(a) = \log(1 - e^a) =
  \begin{cases}
    \log(-\text{expm1}(a)) & \text{if } a > -\ln 2, \\
    \text{log1p}(-e^a)     & \text{if } a \leq -\ln 2.
  \end{cases}
\]
The case split avoids catastrophic cancellation.

% ====================================================================
\section{Information Gain}
\label{sec:ig}
% ====================================================================

\subsection{Definition}

Let $H_t$ denote the entropy of the belief at time $t$ (precise definition
depends on the method; see below).  The \emph{realized information gain}
at step $t$ is
\[
  \mathrm{IG}_t \;=\; H_t - H_{t+1}.
\]

\subsection{Telescoping and Upper Bound}

\begin{proposition}[Cumulative IG telescopes]
\label{prop:telescope}
The cumulative realized information gain over $T$ steps satisfies
\[
  \sum_{t=0}^{T-1} \mathrm{IG}_t
  = \sum_{t=0}^{T-1} (H_t - H_{t+1})
  = H_0 - H_T.
\]
\end{proposition}

\begin{proof}
Immediate by telescoping.  The sum collapses because $H_{t+1}$ in term $t$
cancels $H_{t+1}$ in term $t{+}1$.
\end{proof}

\begin{proposition}[Universal upper bound]
\label{prop:upperbound}
Since entropy is non-negative ($H_T \geq 0$), we have
\[
  \sum_{t=0}^{T-1} \mathrm{IG}_t = H_0 - H_T \;\leq\; H_0.
\]
When the prior is independent Bernoulli,
\[
  H_0 = \sum_{e=1}^{m} h(p_e^{(0)}).
\]
This bound is tight: it is achieved when all edges are resolved
($H_T = 0$, i.e.\ the posterior is deterministic).
\end{proposition}

\subsection{Expected Information Gain}

Before observing the outcome of query $\mathbf{q}$, we can compute the
\emph{expected} IG.  Let $p_{\text{col}} = P(r=1 \mid \mathbf{q})$ be
the predictive collision probability under the current posterior.  The
expected information gain is bounded by
\[
  \mathbb{E}[\mathrm{IG}] \;\leq\; h(p_{\text{col}}).
\]
Equality holds when the query fully resolves the collision event (i.e.,
the observation is a sufficient statistic for the queried edges).

% ====================================================================
\section{Method 1: Approximate (Factored / Mean-Field) Update}
\label{sec:approximate}
% ====================================================================

\subsection{State Representation}

The approximate method maintains a factored belief:
\[
  Q_t(\mathbf{z})
  = \prod_{e=1}^{m} p_e^{(t)\,z_e}\,\bigl(1-p_e^{(t)}\bigr)^{1-z_e},
\]
where the per-edge beliefs are stored internally as log-odds
$\ell_e^{(t)} = \logodds(p_e^{(t)})$.  Additionally, a set
$\mathcal{S}_t \subseteq E$ of \emph{confirmed-safe} edges is maintained;
for $e \in \mathcal{S}_t$, $\ell_e = -\infty$ (i.e.\ $p_e = 0$)
permanently.

\subsection{Safe Observation Update ($r = 0$)}

When $r = 0$, all queried edges are known hazard-free with certainty.
This is an \emph{exact} update:
\begin{align}
  \forall\, e \text{ with } q_e = 1:\quad
  &p_e^{(t+1)} = 0, \quad
   \ell_e^{(t+1)} = -\infty, \quad
   \mathcal{S}_{t+1} = \mathcal{S}_t \cup \{e : q_e = 1\}. \label{eq:approx-safe}
\end{align}
Non-queried edges are unchanged: $p_e^{(t+1)} = p_e^{(t)}$.

\begin{remark}
The safe update is exact because the observation $r=0$ factorizes:
$P(r{=}0 \mid \mathbf{Z}) = \prod_{e:\,q_e=1} (1 - Z_e)$, so conditioning
preserves independence.  No approximation is needed.
\end{remark}

\subsection{Collision Observation Update ($r = 1$)}
\label{sec:approx-collision}

When $r = 1$, we know $\bigvee_{e:\,q_e=1} Z_e = 1$, i.e.\ at least one
queried edge is hazardous.  Let
\[
  Q = \{e : q_e = 1,\; e \notin \mathcal{S}_t\}
\]
be the set of queried, unresolved edges.  Define
\[
  P_{\text{safe}}
  = P\Bigl(\bigwedge_{e \in Q} Z_e = 0\Bigr)
  = \prod_{e \in Q} (1 - p_e^{(t)}).
\]
The probability of collision is $P_{\text{col}} = 1 - P_{\text{safe}}$,
and the scaling factor is
\begin{equation}
  \alpha = \frac{1}{P_{\text{col}}} = \frac{1}{1 - P_{\text{safe}}}.
  \label{eq:alpha}
\end{equation}

In log-space:
\begin{align}
  \log P_{\text{safe}} &= \sum_{e \in Q} \log(1 - p_e) = -\sum_{e \in Q} \softplus(\ell_e), \label{eq:log-psafe} \\
  \log \alpha &= -\texttt{log1mexp}\bigl(\log P_{\text{safe}}\bigr). \label{eq:log-alpha}
\end{align}

Each queried edge is then updated:
\begin{equation}
  p_e^{(t+1)} = \alpha \cdot p_e^{(t)}, \qquad e \in Q.
  \label{eq:alpha-scale}
\end{equation}

In log-odds form:
\begin{align}
  \log(\alpha\, p_e) &= \log\alpha + \log p_e
  = \log\alpha - \softplus(-\ell_e), \label{eq:log-alpha-p} \\[4pt]
  \ell_e^{(t+1)} &= \log\frac{\alpha\, p_e}{1 - \alpha\, p_e}
  = \log(\alpha\, p_e) - \texttt{log1mexp}\bigl(\log(\alpha\, p_e)\bigr).
  \label{eq:new-logodds}
\end{align}
If $\log(\alpha\, p_e) \geq 0$ (i.e.\ $\alpha\, p_e \geq 1$), we set
$\ell_e^{(t+1)} = +\infty$ (saturation to $p_e = 1$).

Non-queried edges are unchanged: $\ell_e^{(t+1)} = \ell_e^{(t)}$.

\subsubsection{Derivation of the $\alpha$-Scaling Rule}

The exact single-edge posterior under a collision observation is:
\begin{equation}
  P(Z_e = 1 \mid r{=}1)
  = \frac{P(r{=}1 \mid Z_e = 1)\, P(Z_e = 1)}{P(r{=}1)}.
  \label{eq:bayes-single}
\end{equation}
Since $r = 1$ is guaranteed when $Z_e = 1$ (that edge alone causes a
collision), we have $P(r{=}1 \mid Z_e = 1) = 1$.  Thus:
\begin{equation}
  P(Z_e = 1 \mid r{=}1)
  = \frac{p_e}{P(r{=}1)}
  = \frac{p_e}{P_{\text{col}}}
  = \alpha \cdot p_e.
  \label{eq:alpha-derivation}
\end{equation}

However, this derivation assumes the other queried edges are independent
of $Z_e$ in the posterior.  The observation $r = 1$ creates the
constraint $\bigvee_{e \in Q} Z_e = 1$, which introduces
\emph{negative correlations}: knowing one edge is hazardous makes it
less likely that the others are.  The exact single-edge marginal is
\begin{equation}
  P(Z_e = 1 \mid r{=}1)
  = \frac{p_e \cdot 1 + (1-p_e)\cdot P\bigl(\bigvee_{e' \in Q\setminus\{e\}} Z_{e'} = 1\bigr)}
         {P_{\text{col}}}
  \neq \alpha\, p_e \;\text{in general.}
\end{equation}

Wait---let us be more careful.  We have:
\begin{align}
  P(Z_e = 1 \mid r{=}1)
  &= \frac{P(r{=}1 \mid Z_e = 1)\, p_e}{P_{\text{col}}} \notag \\
  &= \frac{1 \cdot p_e}{P_{\text{col}}} = \alpha\, p_e.
  \label{eq:marginal-exact}
\end{align}
This is actually \emph{exact for the marginal} of each edge individually!
The $\alpha$-scaling correctly computes $P(Z_e = 1 \mid r = 1)$.

\begin{proposition}[Marginal correctness of $\alpha$-scaling]
\label{prop:marginal-correct}
Under the prior independence assumption, for any single queried edge $e$:
\[
  P(Z_e = 1 \mid r{=}1) = \frac{p_e}{P_{\text{col}}} = \alpha\, p_e.
\]
\end{proposition}

\begin{proof}
$P(r{=}1 \mid Z_e = 1) = 1$ because edge $e$ is on the query and $Z_e = 1$
guarantees a collision.  Apply Bayes' rule directly.
\end{proof}

\begin{remark}[Where the approximation lies]
The updated marginals $\{p_e^{(t+1)}\}$ are individually correct, but the
approximate method \emph{re-factors} the belief as a product
$Q_{t+1}(\mathbf{z}) = \prod_e \Ber(p_e^{(t+1)})$.  The true posterior
$P(\mathbf{Z} \mid r{=}1)$ is \emph{not} a product distribution---the
edges are negatively correlated under the constraint
$\bigvee Z_e = 1$.  Specifically, the joint satisfies:
\[
  P(\mathbf{Z} = \mathbf{z} \mid r{=}1)
  = \frac{\mathbf{1}[\mathbf{z} \cdot \mathbf{q} > 0] \cdot P(\mathbf{z})}{P_{\text{col}}}
  \neq \prod_e p_e^{(t+1)\,z_e}(1-p_e^{(t+1)})^{1-z_e}.
\]
The approximate method discards these correlations at each step.
On subsequent queries, these lost correlations compound---the
approximate marginals may diverge from the true marginals because
the ``prior'' for the next step is the refactored product, not
the true correlated posterior.
\end{remark}

\subsection{Entropy and Information Gain}

Since $Q_t$ is a product distribution, its entropy decomposes:
\begin{equation}
  H(Q_t) = \sum_{e=1}^{m} h(p_e^{(t)}) = \sum_{e=1}^m h\bigl(\sigmoid(\ell_e^{(t)})\bigr).
  \label{eq:approx-entropy}
\end{equation}
Using the log-odds form \eqref{eq:hlogodds}, this is computed without
converting to probabilities.  The realized IG is
\[
  \mathrm{IG}_t^{\text{approx}} = H(Q_t) - H(Q_{t+1}).
\]

\subsubsection{When Can Approximate IG Be Negative?}

For the safe case ($r = 0$), confirmed-safe edges move from $p_e > 0$ to
$p_e = 0$, strictly reducing $h(p_e)$.  So $\mathrm{IG}_t \geq 0$ always
when $r = 0$.

For the collision case ($r = 1$), the $\alpha$-scaling increases each
queried edge's probability.  If $p_e^{(t)} < 0.5$ and
$\alpha\, p_e^{(t)}$ lands closer to $0.5$, then $h(p_e)$ \emph{increases}.
Since $h$ is maximized at $p = 0.5$, pushing probabilities toward $0.5$
increases marginal entropy, and
$\mathrm{IG}_t^{\text{approx}} < 0$ can occur.

\begin{lemma}[Condition for negative approximate IG on collision]
$\mathrm{IG}_t^{\text{approx}} < 0$ can occur when $r^{(t)} = 1$ if and
only if the total increase in entropy from edges pushed toward $0.5$
exceeds the decrease (if any) from edges pushed past $0.5$:
\[
  \sum_{e \in Q} \bigl[h(\alpha\, p_e) - h(p_e)\bigr] > 0.
\]
\end{lemma}

This is a known limitation of any factored approximation applied to
disjunctive constraints.

\subsection{Lifted Product Joint}
\label{sec:lifted-joint}

For apples-to-apples comparison with the exact method at the
joint-distribution level, we define the \emph{lifted product joint}:
\[
  Q_t(\mathbf{z}) = \prod_{e=1}^{m} p_e^{(t)\,z_e}\,(1-p_e^{(t)})^{1-z_e},
  \qquad \mathbf{z} \in \{0,1\}^m.
\]
This is stored as probabilities over $2^m$ states in bitmask order
(bit $e$ of the state index $\leftrightarrow$ $z_e$), matching the exact
method's representation.  In log-space:
\begin{equation}
  \log Q_t(\mathbf{z})
  = \sum_{e=1}^{m}
    \begin{cases}
      -\softplus(-\ell_e^{(t)}) & \text{if } z_e = 1, \\[2pt]
      -\softplus(\ell_e^{(t)})  & \text{if } z_e = 0.
    \end{cases}
  \label{eq:lifted-log}
\end{equation}

\begin{proposition}[Entropy of lifted product joint]
\label{prop:lifted-entropy}
For the product distribution $Q_t$:
\[
  H(Q_t) = -\sum_{\mathbf{z} \in \{0,1\}^m} Q_t(\mathbf{z})\,\log_2 Q_t(\mathbf{z})
  = \sum_{e=1}^{m} h(p_e^{(t)}).
\]
\end{proposition}

\begin{proof}
By independence, the joint entropy equals the sum of marginal entropies.
Formally, since $Q_t(\mathbf{z}) = \prod_e Q_t^{(e)}(z_e)$ where
$Q_t^{(e)}$ is the Bernoulli marginal for edge $e$:
\begin{align*}
  H(Q_t)
  &= -\sum_{\mathbf{z}} \prod_e Q_t^{(e)}(z_e) \cdot \sum_{e'} \log_2 Q_t^{(e')}(z_{e'}) \\
  &= -\sum_{e'} \sum_{z_{e'} \in \{0,1\}} Q_t^{(e')}(z_{e'}) \log_2 Q_t^{(e')}(z_{e'})
     \cdot \underbrace{\prod_{e \neq e'} \sum_{z_e} Q_t^{(e)}(z_e)}_{= 1} \\
  &= \sum_{e=1}^{m} h(p_e^{(t)}).  \qedhere
\end{align*}
\end{proof}

% ====================================================================
\section{Method 2: Exact (Full Joint Bitmask) Update}
\label{sec:exact}
% ====================================================================

\subsection{State Representation}

The exact method maintains the full joint posterior over all $2^m$ states,
stored as log-probabilities:
\[
  \lambda_{\mathbf{z}}^{(t)}
  = \log P(\mathbf{Z} = \mathbf{z} \mid \mathcal{D}_t),
  \qquad \mathbf{z} \in \{0,1\}^m.
\]
State $\mathbf{z}$ is identified with integer
$z = \sum_{e=1}^m z_e \cdot 2^{e-1}$, so edge $e$ corresponds to bit
$e{-}1$ of the index.  The constraint $m \leq 62$ ensures the bitmask
fits in a 64-bit integer.

\subsection{Prior Construction}
\label{sec:exact-prior}

Under the independent prior, the log-probability of state $\mathbf{z}$ is:
\begin{equation}
  \log P(\mathbf{z})
  = \sum_{e=1}^{m} \bigl[z_e \log p_e + (1-z_e)\log(1-p_e)\bigr]
  = \sum_{e=1}^{m}
    \begin{cases}
      -\softplus(-\ell_e^{(0)}) & \text{if } z_e = 1, \\[2pt]
      -\softplus(\ell_e^{(0)})  & \text{if } z_e = 0,
    \end{cases}
  \label{eq:exact-prior}
\end{equation}
using identities \eqref{eq:logp}--\eqref{eq:log1mp}.
After construction, the log-posterior is normalized:
$\lambda_{\mathbf{z}} \gets \lambda_{\mathbf{z}} - \logsumexp_{\mathbf{z}'} \lambda_{\mathbf{z}'}$.

\begin{remark}
At the prior, the exact and approximate representations are equivalent:
both encode the same product distribution.  The log-posterior
\eqref{eq:exact-prior} is identical to the lifted product log-joint
\eqref{eq:lifted-log} with $t = 0$.
\end{remark}

\subsection{Observation Update (Bayesian Conditioning)}
\label{sec:exact-update}

Define the query bitmask $M = \sum_{e:\,q_e=1} 2^{e-1}$.

\paragraph{Likelihood.}
The observation likelihood for state $\mathbf{z}$ is:
\begin{equation}
  P(r \mid \mathbf{Z} = \mathbf{z}, \mathbf{q})
  = \begin{cases}
    1 & \text{if } r = 1 \text{ and } (\mathbf{z}\;\mathsf{AND}\; M) \neq 0, \\
    1 & \text{if } r = 0 \text{ and } (\mathbf{z}\;\mathsf{AND}\; M) = 0, \\
    0 & \text{otherwise}.
  \end{cases}
  \label{eq:likelihood}
\end{equation}
In words: collision requires at least one hazardous edge on the path;
safe requires zero.

\paragraph{Posterior update.}
By Bayes' theorem:
\begin{equation}
  P(\mathbf{z} \mid \mathcal{D}_{t+1})
  = \frac{P(r^{(t)} \mid \mathbf{z}, \mathbf{q}^{(t)})\, P(\mathbf{z} \mid \mathcal{D}_t)}
         {P(r^{(t)} \mid \mathbf{q}^{(t)}, \mathcal{D}_t)}.
  \label{eq:bayes-exact}
\end{equation}
Since the likelihood \eqref{eq:likelihood} is 0--1 valued, this reduces
to zeroing inconsistent states and renormalizing:

\begin{algorithm}[H]
\caption{Exact observation update}
\begin{algorithmic}[1]
\Require log-posterior $\{\lambda_{\mathbf{z}}\}$, query bitmask $M$, outcome $r$
\For{each state $\mathbf{z} \in \{0,1\}^m$}
  \State $\text{hit} \gets (\mathbf{z}\;\mathsf{AND}\; M) \neq 0$
  \If{$(r = 1 \;\wedge\; \neg\text{hit})$ or $(r = 0 \;\wedge\; \text{hit})$}
    \State $\lambda_{\mathbf{z}} \gets -\infty$ \Comment{$P(\mathbf{z} \mid \mathcal{D}_{t+1}) = 0$}
  \EndIf
\EndFor
\State $c \gets \logsumexp_{\mathbf{z}} \lambda_{\mathbf{z}}$ \Comment{$\log P(r \mid \mathbf{q}, \mathcal{D}_t)$}
\State $\lambda_{\mathbf{z}} \gets \lambda_{\mathbf{z}} - c$ for all $\mathbf{z}$
  \Comment{normalize}
\end{algorithmic}
\end{algorithm}

The normalizing constant $c$ equals $\log P(r \mid \mathbf{q}, \mathcal{D}_t)$,
the log-evidence of the observation.  If $c = -\infty$, the observations
are inconsistent and an error is raised.

\paragraph{Complexity.}
Each update iterates over all $2^m$ states, so the cost is $O(2^m)$.

\subsection{Derived Quantities}

\subsubsection{Joint Entropy}
\begin{equation}
  H_{\text{joint}}^{(t)}
  = -\sum_{\mathbf{z} \in \{0,1\}^m} P(\mathbf{z} \mid \mathcal{D}_t)\,\log_2 P(\mathbf{z} \mid \mathcal{D}_t)
  = -\frac{1}{\ln 2}\sum_{\mathbf{z}} e^{\lambda_{\mathbf{z}}^{(t)}} \cdot \lambda_{\mathbf{z}}^{(t)}.
  \label{eq:joint-entropy}
\end{equation}
Terms with $\lambda_{\mathbf{z}} = -\infty$ contribute zero (by the
convention $0 \cdot (-\infty) = 0$).

\subsubsection{Marginal Probabilities}
\begin{equation}
  P(Z_e = 1 \mid \mathcal{D}_t)
  = \sum_{\mathbf{z}:\, z_e = 1} P(\mathbf{z} \mid \mathcal{D}_t)
  = \exp\!\Bigl(\logsumexp_{\mathbf{z}:\, z_e = 1} \lambda_{\mathbf{z}}^{(t)}\Bigr).
  \label{eq:exact-marginal}
\end{equation}

\subsubsection{Marginal Entropy}
\begin{equation}
  H_{\text{marginal}}^{(t)} = \sum_{e=1}^{m} h\!\bigl(P(Z_e = 1 \mid \mathcal{D}_t)\bigr).
  \label{eq:marginal-entropy}
\end{equation}

\subsubsection{Predictive Collision Probability}
For a candidate query $\mathbf{q}$ with bitmask $M$:
\begin{equation}
  P(r{=}1 \mid \mathbf{q}, \mathcal{D}_t)
  = \sum_{\mathbf{z}:\, (\mathbf{z}\,\mathsf{AND}\, M) \neq 0} P(\mathbf{z} \mid \mathcal{D}_t)
  = \exp\!\Bigl(\logsumexp_{\mathbf{z}:\, (\mathbf{z}\,\mathsf{AND}\, M) \neq 0} \lambda_{\mathbf{z}}^{(t)}\Bigr).
  \label{eq:pred-collision}
\end{equation}

\subsubsection{Information Gain}

The exact method supports two entropy measures:
\begin{equation}
  \mathrm{IG}_t^{\text{exact}} =
  \begin{cases}
    H_{\text{joint}}^{(t)} - H_{\text{joint}}^{(t+1)} & \text{(joint mode)}, \\
    H_{\text{marginal}}^{(t)} - H_{\text{marginal}}^{(t+1)} & \text{(marginal mode)}.
  \end{cases}
  \label{eq:exact-ig}
\end{equation}
Joint IG is always non-negative (conditioning never increases Shannon
entropy).  Marginal IG can in principle be negative (marginal entropies
can increase even when the joint entropy decreases, though this is rare
for the exact method).

% ====================================================================
\section{Relationship Between Joint and Marginal Entropy}
\label{sec:joint-vs-marginal}
% ====================================================================

\begin{proposition}[Subadditivity of entropy]
\label{prop:subadditivity}
For any joint distribution $P$ on $(Z_1, \dots, Z_m)$:
\[
  H(Z_1, \dots, Z_m) \;\leq\; \sum_{e=1}^{m} H(Z_e),
\]
with equality if and only if $Z_1, \dots, Z_m$ are mutually independent.
\end{proposition}

\begin{proof}
This is a standard result from information theory.  By the chain rule:
\[
  H(Z_1, \dots, Z_m) = \sum_{e=1}^{m} H(Z_e \mid Z_1, \dots, Z_{e-1})
  \leq \sum_{e=1}^{m} H(Z_e),
\]
since conditioning reduces entropy: $H(X \mid Y) \leq H(X)$.  Equality
holds iff each $H(Z_e \mid Z_1, \dots, Z_{e-1}) = H(Z_e)$, which is the
definition of mutual independence.
\end{proof}

\paragraph{Consequences.}
\begin{enumerate}
  \item \textbf{At the prior}: the edges are independent, so
    $H_{\text{joint}}^{(0)} = H_{\text{marginal}}^{(0)} = \sum_e h(p_e^{(0)})$.
  \item \textbf{After safe observations}: safe observations preserve
    independence (the likelihood factorizes), so
    $H_{\text{joint}} = H_{\text{marginal}}$ still holds.
  \item \textbf{After collision observations}: the constraint
    $\bigvee_{e \in Q} Z_e = 1$ introduces correlations.  Then
    $H_{\text{joint}} < H_{\text{marginal}}$.  The gap
    $H_{\text{marginal}} - H_{\text{joint}}$ is the total correlation
    (multi-information):
    \begin{equation}
      \mathrm{TC}(\mathbf{Z} \mid \mathcal{D}_t)
      = \sum_{e=1}^m H(Z_e \mid \mathcal{D}_t) - H(\mathbf{Z} \mid \mathcal{D}_t)
      = D_{\mathrm{KL}}\!\bigl(P(\mathbf{Z} \mid \mathcal{D}_t) \;\|\; Q_t\bigr),
      \label{eq:total-correlation}
    \end{equation}
    where $Q_t = \prod_e P(Z_e \mid \mathcal{D}_t)$ is the
    marginal-product approximation.
\end{enumerate}

% ====================================================================
\section{Comparison of Methods}
\label{sec:comparison}
% ====================================================================

\begin{table}[h]
\centering
\begin{tabular}{@{}lcc@{}}
\toprule
                        & \textbf{Approximate}             & \textbf{Exact} \\
\midrule
State size              & $O(m)$ log-odds                  & $O(2^m)$ log-probs \\
Update cost             & $O(m)$ per step                  & $O(2^m)$ per step \\
Handles correlations    & No (factored)                    & Yes (full joint) \\
Safe observation        & Exact                            & Exact \\
Collision marginals     & Exact (for one step)             & Exact \\
Collision joint         & Approximate (refactored)         & Exact \\
Joint entropy           & $\sum_e h(p_e)$                  & $-\sum_{\mathbf{z}} P(\mathbf{z})\log_2 P(\mathbf{z})$ \\
IG always $\geq 0$?    & No (collision can be negative)   & Yes (joint), mostly (marginal) \\
Edge limit              & None                             & $m \leq 62$ (bitmask) \\
\bottomrule
\end{tabular}
\caption{Summary comparison of the two update methods.}
\label{tab:comparison}
\end{table}

\subsection{Entropy Ordering After Observations}

After $T$ observations, the various entropy measures satisfy:
\begin{equation}
  H_{\text{exact joint}}^{(T)}
  \;\leq\;
  H_{\text{exact marginal}}^{(T)}
  = H_{\text{approx}}^{(T)} + \Delta_T,
  \label{eq:entropy-ordering}
\end{equation}
where $\Delta_T$ accounts for the difference in marginals between the
exact and approximate methods.  After only safe observations,
$\Delta_T = 0$ and all three are equal.  After collision observations,
the approximate marginals may diverge from exact marginals, and
the exact joint is strictly less than both marginal sums.

\subsection{Why the Approximate Method Loses Information}

Consider two successive observations:
\begin{enumerate}
  \item Query $\mathbf{q}^{(0)}$, observe $r^{(0)} = 1$ (collision).
        The true posterior $P(\mathbf{Z} \mid r^{(0)} = 1)$ has
        correlations; the approximate method replaces it with
        $Q_1 = \prod_e \Ber(\alpha\, p_e)$.
  \item Query $\mathbf{q}^{(1)}$.  The exact method conditions
        $P(\mathbf{Z} \mid r^{(0)} = 1)$ on the new observation.
        The approximate method conditions $Q_1$ instead---a
        different distribution.
\end{enumerate}
The KL divergence $D_{\mathrm{KL}}(P \| Q_1)$ accumulated at step 1
means the approximate ``prior'' for step 2 is wrong.  This error
compounds across steps.  In the exact method, no information is
discarded.

\subsection{Universal Bounds}

For both methods, the cumulative IG satisfies:
\begin{equation}
  0 \;\leq\; \sum_{t=0}^{T-1} \mathrm{IG}_t^{(\cdot)} = H_0^{(\cdot)} - H_T^{(\cdot)}
  \;\leq\; H_0^{(\cdot)}.
\end{equation}
When the prior is independent Bernoulli with per-edge probability $p$:
\[
  H_0 = m \cdot h(p).
\]
For $p = 0.5$: $H_0 = m$ bits (one bit of uncertainty per edge).

\end{document}
